% Standalone problem document, generated from the adjacent README.md.
% Compile with XeLaTeX. Mathematical content is maintained in Markdown.
\documentclass[11pt,a4paper]{article}
\usepackage[fontsize=10.5pt]{scrextend}
\usepackage[margin=22mm,headheight=15pt,headsep=9mm,footskip=11mm]{geometry}
\usepackage{fontspec}
\setmainfont{texgyrepagella}[Extension=.otf,UprightFont=*-regular,BoldFont=*-bold,ItalicFont=*-italic,BoldItalicFont=*-bolditalic]
\setsansfont{texgyreheros}[Extension=.otf,UprightFont=*-regular,BoldFont=*-bold,ItalicFont=*-italic,BoldItalicFont=*-bolditalic]
\setmonofont{texgyrecursor}[Extension=.otf,UprightFont=*-regular,BoldFont=*-bold,ItalicFont=*-italic,BoldItalicFont=*-bolditalic,Scale=MatchLowercase]
\usepackage{amsmath,amssymb}
\usepackage{unicode-math}
\setmathfont{texgyrepagella-math.otf}
\usepackage{xcolor,microtype,enumitem,needspace,fancyhdr}
\usepackage{bookmark}
\definecolor{ink}{HTML}{17344B}
\definecolor{accent}{HTML}{197D80}
\definecolor{muted}{HTML}{596773}
\hypersetup{colorlinks=true,linkcolor=accent,urlcolor=accent,pdftitle={RA-11 - Optimal
trace-estimation complexity using Kronecker matrix-vector
queries},pdfauthor={Open Problems in Numerical Linear Algebra}}
\urlstyle{same}
\setlength{\parindent}{0pt}
\setlength{\parskip}{5pt plus 1pt minus 1pt}
\linespread{1.04}
\setlength{\emergencystretch}{3em}
\widowpenalty=10000
\clubpenalty=10000
\displaywidowpenalty=10000
\allowdisplaybreaks[1]
\setlist{nosep,leftmargin=1.5em}
\providecommand{\tightlist}{\setlength{\itemsep}{0pt}\setlength{\parskip}{0pt}}
\providecommand{\pandocbounded}[1]{#1}
\pagestyle{fancy}
\fancyhf{}
\fancyhead[L]{\sffamily\footnotesize\color{muted}OPEN PROBLEMS IN NUMERICAL LINEAR ALGEBRA}
\fancyhead[R]{\sffamily\bfseries\footnotesize\color{accent}RA-11}
\fancyfoot[L]{\parbox[b]{0.9\textwidth}{\sffamily\fontsize{8}{10}\selectfont\color{muted}Editorial ratings; a bounded search cannot certify that a problem remains unsolved.\\Literature check: 2026-09-08}}
\fancyfoot[R]{\sffamily\footnotesize\color{muted}\thepage}
\renewcommand{\headrulewidth}{0.3pt}
\renewcommand{\headrule}{\hbox to\headwidth{\color{accent}\leaders\hrule height \headrulewidth\hfill}}
\makeatletter
\renewcommand{\section}{\Needspace{5\baselineskip}\@startsection{section}{1}{0pt}{12pt plus 2pt minus 2pt}{4pt}{\sffamily\bfseries\large\color{ink}}}
\renewcommand{\subsection}{\Needspace{4\baselineskip}\@startsection{subsection}{2}{0pt}{9pt plus 1pt minus 1pt}{3pt}{\sffamily\bfseries\normalsize\color{ink}}}
\makeatother
\setcounter{secnumdepth}{0}
\begin{document}
{\sffamily\small\color{muted}Randomized and low-rank approximation\par}
\vspace{3pt}
{\raggedright\hyphenpenalty=10000\exhyphenpenalty=10000\sffamily\bfseries\fontsize{21}{24}\selectfont\color{ink}Optimal
trace-estimation complexity using Kronecker matrix-vector queries\par}
\vspace{7pt}
{\sffamily\small\textbf{Difficulty:} extreme\quad\textbf{Importance:} interesting
to the community\par}
\vspace{4pt}
{\color{accent}\hrule height 0.8pt}
\vspace{6pt}
\textbf{Topic:} Structured randomized trace estimation.

Let \(n,q\ge2\) be integers and \(0<\varepsilon<1/2\). An unknown real
symmetric PSD matrix \(M\in\mathbb R^{n^q\times n^q}\) is accessible
only through an exact oracle which, on input
\(v_1,\ldots,v_q\in\mathbb R^n\), returns

\[
M(v_1\otimes\cdots\otimes v_q)\in\mathbb R^{n^q}.
\]

Let \(Q(n,q,\varepsilon)\) be the smallest integer \(t\) for which a
randomized algorithm can, for every such \(M\), make at most \(t\)
oracle calls and output \(\widehat t\in\mathbb R\) satisfying

\[
\Pr\bigl\{|\widehat t-\operatorname{tr}(M)|\le\varepsilon\operatorname{tr}(M)\bigr\}\ge\frac23.
\]

Queries may depend on previous responses and internal randomness. The
algorithm may perform arbitrary finite exact arithmetic between calls;
only oracle calls are counted. The inputs \(n,q,\varepsilon\) are known.
Neither the query vectors nor their concatenation are required to have
bounded condition number, and there is no restriction to Hutchinson
estimators.

Determine \(Q(n,q,\varepsilon)\) up to universal constant factors,
simultaneously in all three parameters.

This is a precise query-complexity formulation of the source's request
for tight trace-estimation bounds. It would quantify the cost of
exploiting rank-one tensor queries when ordinary matrix access is
unavailable.

\subsection{References}\label{references}

\begin{enumerate}
\def\labelenumi{\arabic{enumi}.}
\tightlist
\item
  R. A. Meyer, W. Swartworth, and D. P. Woodruff,
  \href{https://arxiv.org/html/2502.08029v2}{Understanding the Kronecker
  Matrix-Vector Complexity of Linear Algebra}, ICML 2025, PMLR 267,
  43909--43933. Definition 1 specifies the oracle; §6, final sentence,
  asks for tight trace-estimation bounds. Theorem 7 bounds a conditioned
  scalar quadratic-form oracle; its following discussion distinguishes
  full-vector responses.
  \href{https://proceedings.mlr.press/v267/meyer25a.html}{Published
  record}.
\item
  R. A. Meyer and H. Avron,
  \href{https://arxiv.org/abs/2309.04952}{Hutchinson's Estimator is Bad
  at Kronecker-Trace-Estimation}, \emph{SIAM Journal on Matrix Analysis
  and Applications} 47 (2026), 353--387, abstract and trace-estimator
  bounds. These tight rates concern the specified estimator and query
  distributions.
\end{enumerate}

\subsection{Status check --- 2026-09-08}\label{status-check-2026-09-08}

Checked the first paper's latest listed arXiv v2 (2025-02-13), §6, and
its ICML publication record; checked the second paper's latest listed v2
(2025-01-31) and 2026 journal record. Searches for ``Kronecker trace
estimation tight 2026'' and the exact first title found no solution for
unrestricted adaptive algorithms. The available trace lower bound
assumes both scalar quadratic-form queries and bounded query
conditioning; it does not establish a lower bound for the unrestricted
full-vector oracle defined here. The 2026 Hutchinson article establishes
estimator-specific rates and therefore does not determine \(Q\). The
minimax definition and explicit parameter range are editorial
formalization of a source-stated complexity question, not a quoted
formula from the paper.
\end{document}
